\section*{अध्याय \ref{ch:b3:x-ray-diffraction} --- क्रिस्टलिकी और एक्स-किरण विवर्तन}

\begin{solution}{exo:b3:x-ray-diffraction:1}
$d = a/\sqrt N$: $d_{111} = \qty{325.66}{pm}$, $d_{200} = \qty{282.03}{pm}$, $d_{220} =
\qty{199.42}{pm}$।
\end{solution}

\begin{solution}{exo:b3:x-ray-diffraction:2}
FCC: पहले परावर्तन (111) और (200) हैं। $\sin\theta = \lambda\sqrt N/2a$: $2\theta =
\ang{43.32}$ और \ang{50.45}।
\end{solution}

\begin{solution}{exo:b3:x-ray-diffraction:3}
P: सातों। I ($h + k + l$ सम): 110, 200, 211, 220। F (अमिश्रित): 111, 200, 220।
\end{solution}

\begin{solution}{exo:b3:x-ray-diffraction:4}
$\mathbf a^* = \mathbf e_x/a$, $\mathbf b^* = \mathbf e_y/b$: भुजाओं $1/a$
और $1/b$ वाला एक आयताकार जालक, जिसमें प्रत्यक्ष जालक की लंबी भुजा छोटी बन जाती है।
\end{solution}

\begin{solution}{exo:b3:x-ray-diffraction:5}
$\sin^2\theta$ के अनुपात $1 : 2 : 3 : 4 : 5$, अर्थात् $N = 2, 4, 6, 8, 10$: अंतःकेंद्रित (110,
200, 211, 220, 310)। पहली रेखा से $d_{110} = \lambda/2\sin(\ang{20.135}) =
\qty{223.77}{pm}$ और $a = d\sqrt2 = \qty{316.5}{pm}$: टंगस्टन।
\end{solution}

\begin{solution}{exo:b3:x-ray-diffraction:6}
$Z = \rho N_Aa^3/M = 1.99 \times 6.022\times10^{23} \times (6.2929\times10^{-8})^3/74.6 = 4.0$।
\end{solution}

\begin{solution}{exo:b3:x-ray-diffraction:7}
$\beta = \qty{0.40}{\degree} = \qty{6.98e-3}{rad}$, $\cos\theta = \cos\ang{15.85} = 0.962$:
$L = 0.9 \times 0.15406/(6.98\times10^{-3} \times 0.962) = \qty{21}{nm}$।
\end{solution}

\begin{solution}{exo:b3:x-ray-diffraction:8}
$F_{hkl} = f_{\ce{Cs+}} + f_{\ce{Cl-}}(-1)^{h+k+l}$: $F_{100} = 54 - 18 = 36$, $F_{110} = 72$:
तीव्रताएँ 1 : 4 के अनुपात में। केंद्र पर कोनों वाले आयन से भिन्न आयन
है: जालक आद्य है (अंतःकेंद्रित जालक के लिए दोनों बिंदुओं पर सर्वसम
सामग्री चाहिए), और 100 उपस्थित है।
\end{solution}

\begin{solution}{exo:b3:x-ray-diffraction:9}
हर परमाणु की एक प्रति $(\frac12,\frac12,0)$ से विस्थापित है, जो $F$ को $1 +
\eu^{\iu\pi(h+k)} = 1 + (-1)^{h+k}$ से गुणा करती है, जो $h + k$ विषम होने पर शून्य है।
\end{solution}

\begin{solution}{exo:b3:x-ray-diffraction:10}
अर्धक्रिस्टलों में आवर्तिता के बिना दीर्घ-परास क्रम होता है: उनका कोई
जालक नहीं होता, अतः प्रतिबंध, जो जालकों से संबंधित है, उन पर लागू नहीं होता।
उनके तीक्ष्ण विवर्तन धब्बों के कारण क्रिस्टल की परिभाषा ऐसे किसी भी ठोस के रूप में की गई
जिसका विवर्तन प्रतिरूप मूलतः विविक्त हो, आवर्ती हो या नहीं।
\end{solution}

\begin{solution}{exo:b3:x-ray-diffraction:11}
$c$ सर्पण $(x, y, z)$ को $(x, \frac12 - y, \frac12 + z)$ पर प्रतिचित्रित करता है। $k = 0$ के लिए दोनों
कला गुणक $\eu^{2\pi\iu(hx + lz)}$ और $\eu^{2\pi\iu(hx + lz + l/2)} = (-1)^l
\eu^{2\pi\iu(hx+lz)}$ हैं: उनका योग विषम $l$ के लिए लुप्त होता है।
\end{solution}

\begin{solution}{exo:b3:x-ray-diffraction:12}
प्रतिलोमन, दर्पण या सर्पण किसी अणु को उसके दर्पण-प्रतिबिंब में बदल देता है, अतः
इनमें से किसी को रखने वाले क्रिस्टल में दोनों प्रतिबिंबरूपी होते। शुद्ध प्रतिबिंबरूपी
केवल उन 65 \omterm{def:b3:x-ray-diffraction:space-group}{आकाशी समूहों} में क्रिस्टलित होता है जो घूर्णनों, \omterm{def:b3:x-ray-diffraction:space-group}{पेंच अक्षों}
और स्थानांतरणों से बने हैं (ज़ोन्के समूह), जैसे \textit{P}$2_12_12_1$। फ़्रीडेल के नियम के अनुसार
एक प्रतिबिंबरूपी का प्रतिरूप दूसरे के प्रतिरूप के बराबर है; निरपेक्ष
विन्यास भारी परमाणुओं के असामान्य प्रकीर्णन से उत्पन्न छोटे विचलनों से
प्राप्त किया जाता है।
\end{solution}

\begin{solution}{pb:b3:x-ray-diffraction:1}
\textbf{1.} $\theta = \ang{14.171}$, $d = \lambda/2\sin\theta = \qty{314.64}{pm}$।
\textbf{2.} 314.64, 222.49, 181.66, 157.32, 140.71, \qty{128.45}{pm}।
\textbf{3.} 1.000, 2.000, 3.000, 4.000, 5.000, 6.000।
\textbf{4.} 1, 2, 3, 4, 5, 6।
\textbf{5.} हाँ, $a' = \qty{314.6}{pm}$ वाली आद्य घनीय कोष्ठिका के (100), (110), (111), (200), (210), (211)
के रूप में।
\textbf{6.} $N = 7$ (और 15, 23\dots), जो तीन वर्गों का योग नहीं है; पहला
अंतर आठवीं रेखा पर आता। छह रेखाएँ कोर $a'$ वाली P
कोष्ठिका और कोर $2a'$ वाली F कोष्ठिका में भेद नहीं कर सकतीं।
\textbf{7.} $N = 4, 8, 12, 16, 20, 24$: (200), (220), (222), (400), (420), (422)।
\textbf{8.} $a = 2d_{200} = \qty{629.3}{pm}$।
\textbf{9.} \ce{KCl} (\qty{629.29}{pm})।
\textbf{10.} सभी-विषम सूचकांकों के लिए $F = 4[f(\ce{K+}) - f(\ce{Cl-})]$, और दोनों आयनों के
18-18 इलेक्ट्रॉन हैं: आयाम निरस्त हो जाते हैं।
\textbf{11.} दोनों के लिए हाँ: \ce{Na+} (10) और \ce{Cl-} (18), \ce{K+} (18) और \ce{Br-} (36)
भिन्न रूप से प्रकीर्णन करते हैं; \ce{NaCl} की (111) \ang{27.36} पर होती, \ce{KBr} की \ang{23.38} पर,
दोनों अभिलिखित परास के भीतर।
\textbf{12.} एक्स-किरणें \ce{K+} और \ce{Cl-} को लगभग सर्वसम देखती हैं; सैंधव-लवण व्यवस्था में
सर्वसम आयन कोर $a/2$ वाली एक सरल घनीय व्यूह बनाते हैं, एक आभासी कोष्ठिका।
\textbf{13.} चार \ce{KCl}।
\textbf{14.} $\rho = 4 \times 74.6/(6.022\times10^{23} \times (6.293\times10^{-8})^3) =
\qty{1.99}{g/cm^3}$।
\textbf{15.} छह और छह (दूसरे आयन के अष्टफलक)।
\textbf{16.} $a/2 = \qty{314.6}{pm}$।
\textbf{17.} (440), $N = 32$, \ang{87.65} पर (सभी-विषम (333)/(511) अनुपस्थित होने से)।
\textbf{18.} तीव्रताएँ क्रिस्टलकों के वरीय अभिविन्यास,
अवशोषण, प्रकीर्णन गुणकों के घटने और तापीय गति के साथ बदलती हैं; स्थितियाँ
केवल जालक पर निर्भर करती हैं।
\textbf{19.} $d = \lambda/2\sin\theta$ से $\Delta d/d = -\cot\theta\,\Delta\theta$: वही कोणीय त्रुटि
बड़े $\theta$ पर छोटी आपेक्षिक त्रुटि देती है।
\textbf{20.} $d_{420} = \lambda/2\sin\ang{33.190} = \qty{140.71}{pm}$, $a = d\sqrt{20} =
\qty{629.3}{pm}$।
\textbf{21.} $L = 0.9 \times 0.15406/(4.36\times10^{-3} \times \cos\ang{33.19}) = \qty{38}{nm}$।
\textbf{22.} नहीं: $\beta = \ang{0.0095}$, जो किसी सामान्य उपकरणीय चौड़ाई से बहुत नीचे है।
\textbf{23.} \textbf{$a = \qty{629.3}{pm}$: चूर्ण पोटैशियम क्लोराइड है}, जिसके विषम
परावर्तन समान इलेक्ट्रॉन-संख्या वाले दो आयनों द्वारा मौन कर दिए गए हैं।
\end{solution}
